% from ega0-19-fr.tex
\oldpage[0\textsubscript{IV-1}]{80}

La définition (19.3.1) revient à dire que la propriété suivante a lieu~:
\begin{enumerate}
  \item[(P)] Pour tout idéal ouvert $\mathfrak R$ de la $A$-algèbre $B$ et
  tout $A$-homomorphisme $u':B/\mathfrak R\to C/\mathfrak J$, il y a un
  idéal ouvert $\mathfrak R'\subset\mathfrak R$ de $B$ tel que
  l'homomorphisme
  $B\to B/\mathfrak R\xrightarrow{u'}C/\mathfrak J$ se factorise en
  \[
    B\to B/\mathfrak R'\xrightarrow{v'}C\to C/\mathfrak J,
  \]
  où $v'$ est un $A$-homomorphisme.
\end{enumerate}
En effet, si $u:B\to C/\mathfrak J$ est un $A$-homomorphisme continu, il a
pour noyau $\mathfrak R$ un idéal ouvert de $B$, donc $u$ se factorise en
$B\to B/\mathfrak R\xrightarrow{u'}C/\mathfrak J$, et si (P) est
vérifiée, il suffit de l'appliquer à $u'$ et de prendre pour $v$ la composée
$B\to B/\mathfrak R'\xrightarrow{v'}C$ pour satisfaire aux conditions de
(19.3.1). Inversement, supposons que $B$ soit une $A$-algèbre formellement
lisse~; donnons-nous un idéal ouvert $\mathfrak R$ de $B$ et un
$A$-homomorphisme $u':B/\mathfrak R\to C/\mathfrak J$, et appliquons la
définition (19.3.1) à
$u:B\to B/\mathfrak R\xrightarrow{u'}C/\mathfrak J$~; si
$v:B\to C$ est un $A$-homomorphisme continu tel que $u$ se factorise en
$B\xrightarrow{v}C\to C/\mathfrak J$, l'idéal
$\mathfrak R'=\operatorname{Ker}(v)\cap\mathfrak R$ de $B$ est ouvert et
contenu dans $\mathfrak R$~; par suite $v$ se factorise en
$B\to B/\mathfrak R'\xrightarrow{v'}C$, et $v'$ vérifie bien la condition
(P).

\begin{proposition}[19.3.2]
\label{0.19.3.2-fr}
Soient $A$ un anneau discret, $V$ un $A$-module projectif~; l'algèbre
symétrique $B=\mathbf S_A^\bullet(V)$, munie de la topologie discrète, est une
$A$-algèbre formellement lisse.
\end{proposition}

\begin{proof}
En effet, les notations $C$ et $\mathfrak J$ ayant le même sens que dans
(19.3.1), soit $u:B\to C/\mathfrak J$ un homomorphisme de $A$-algèbres, qui
par restriction à $V=\mathbf S_A^1(V)$, donne un homomorphisme de $A$-modules
$u_1:V\to C/\mathfrak J$~; comme $V$ est projectif, $u_1$ se factorise en
$V\xrightarrow{v_1}C\xrightarrow{\varphi}C/\mathfrak J$, et $v_1$ se
prolonge en un homomorphisme de $A$-algèbres
$v:\mathbf S_A^\bullet(V)\to C$, tel que le composé
$\varphi\circ v$ coïncide avec $u$ dans $V$, donc
$\varphi\circ v=u$, ce qui démontre la proposition.
\end{proof}

\begin{corollary}[19.3.3]
\label{0.19.3.3-fr}
Si $A$ est un anneau discret, toute algèbre de polynômes
$B=A[T_\alpha]_{\alpha\in I}$, munie de la topologie discrète, est une
$A$-algèbre formellement lisse.
\end{corollary}

\begin{proposition}[19.3.4]
\label{0.19.3.4-fr}
Soit $A$ un anneau topologique, et soit
$B=A[[T_\alpha]]_{\alpha\in I}$ un anneau de séries formelles
$\sum_{\lambda\in\mathbb{N}^{(I)}}c_\lambda T^\lambda$ (algèbre large sur
$A$ du monoïde $\mathbb{N}^{(I)}$, identifié en tant que $A$-module au produit
$A^{\mathbb{N}^{(I)}}$). Si l'on munit $B$ de la topologie produit, $B$ est une
$A$-algèbre formellement lisse.
\end{proposition}

Soit $(\mathfrak L_\mu)_{\mu\in M}$ un système fondamental d'idéaux ouverts
dans $A$. Pour toute partie finie $H$ de $I$, tout $\mu\in M$ et tout entier
$n$, désignons par $\mathfrak K_{H,\mu,n}$ le voisinage de $0$ dans $B$ formé
des $(c_\lambda)$ tels que pour tout
$\lambda=(\lambda_\alpha)_{\alpha\in I}$ tel que
$\lambda_\alpha=0$ pour $\alpha\notin H$ et
$\sum_{\alpha\in H}\lambda_\alpha\leq n$, on ait
$c_\lambda\in\mathfrak L_\mu$. On vérifie immédiatement que les
$\mathfrak K_{H,\mu,n}$ forment un système fondamental de voisinages de $0$
dans $B$ et sont des idéaux de $B$, donc la topologie produit est compatible
avec la structure de $A$-algèbre de $B$.

Notons d'abord, avec les mêmes notations, le
\begin{lemma}[19.3.4.1]
\label{0.19.3.4.1-fr}
Soit $E$ une $A$-algèbre discrète.
\begin{enumerate}
  \item[(i)] Si $f:B\to E$ est un $A$-homomorphisme continu, il existe une
  partie finie $H$ de $I$ telle que $f(T_\alpha)=0$ pour $\alpha\notin H$,
  et $f(T_\alpha)$ est nilpotent dans $E$ pour tout $\alpha\in H$.
  \item[(ii)] Inversement, soient $H$ une partie finie de $I$,
  $(z_\alpha)_{\alpha\in H}$ une famille d'éléments nilpotents de $E$.
  Il existe un $A$-homomorphisme continu $g:B\to E$ et un seul tel que
  $g(T_\alpha)=z_\alpha$ pour $\alpha\in H$ et
  $g(T_\alpha)=0$ pour $\alpha\notin H$.
\end{enumerate}
\end{lemma}

